\section{Eigenvalues and Singular Values of a Symmetric Matrix}

Since \(A=A^T\), A has an orthonormal set of vectors and
\[
A=Q\Lambda Q^T
\]
where \(Q\) is orthogonal ($Q^{-1}=Q^T$) and
\[
\Lambda=\operatorname{diag}(\lambda_1,\ldots,\lambda_n)
\]
In general, singular values satisfy
\[
\sigma_i = \sqrt{\lambda_i(A^T A)}
\]
Therefore, as $A$ is symmetric:
\[
A^TA=A^2=Q\Lambda Q^TQ\Lambda Q^T=Q\Lambda^2Q^T
\]
The eigenvalues of \(A^TA\) are \(\lambda_i^2\), so the singular values of
\(A\) are
\[
\sqrt{\lambda_i^2}=|\lambda_i|
\]
Thus
\[
\boxed{\{\sigma_1,\ldots,\sigma_n\}
=\{|\lambda_1|,\ldots,|\lambda_n|\}}
\]
The singular values are the absolute eigenvalues sorted in decreasing order.
% Since the eigenvalues are sorted by value rather than magnitude, it is not
% necessarily true that \(\sigma_i=|\lambda_i|\) for the same index \(i\).

% For an eigenvector \(q_i\) with \(\lambda_i\neq0\), one possible choice of
% right and left singular vectors is
% \[
% v_i=q_i,\qquad u_i=\operatorname{sign}(\lambda_i)q_i
% \]
% because
% \[
% Av_i=|\lambda_i|u_i
% \]
